\section{Conclusion and Future Work} \label{sec:conclusion}
ODDToolkit presents a practical Ontology Driven Design pipeline that keeps documentation, validation artefacts, typed models, and schema-oriented outputs aligned around a shared enriched representation. The environmental reporting project in Section~\ref{sec:usecase} shows that this workflow can be integrated into a development process and can provide useful implementation, documentation, and validation artefacts from one source.

The paper does not yet offer a quantitative evaluation of developer productivity, maintenance cost, or long-term evolution across multiple deployments. However, it demonstrates that good ontology engineering practices can be leveraged to aid in the generation and upkeep of documentation that helps developers understand and work with the underlying data models.

Our approach still depends on explicit project conventions, especially SKOS-based naming and Hydra-based key identification, so its strongest fit remains governed ontologies under active development. It does not yet solve the preservation of customised generated code, populated-database migration, derived indexes or summaries, or partial-update logic. Future work should therefore focus on interoperability with broader schema ecosystems and a broader evaluation across projects.

\paragraph{Supplemental Material.}
ODDToolkit is available on GitHub, including code and website documentation~\cite{oddtoolkit2026}.

\paragraph*{Declaration on Generative AI.}
During the creation of ODDToolkit, the authors used GitHub Copilot to help with the documentation and development. The authors take full responsibility for the content in the repository.